\chapter{设备管理}

\section{I/O 系统与设备分类}
\concept{设备管理的对象与任务}
\begin{points}
  \item 除 CPU 和内存外，大部分硬件设备称为外部设备。
  \item I/O 系统包括 I/O 设备、接口线路、控制部件、通道和管理软件。
  \item I/O 操作是主存与外围设备介质之间的信息传送。
  \item 设备管理任务：响应 I/O 请求、分配 I/O 设备、管理设备可靠性与安全性、提供友好访问接口。
\end{points}
\plain{这题通常先答“设备管理管什么”，再答“为什么需要它”：设备种类多、速度差大、还常被多个进程共享。}
\exam{常考设备管理的对象、任务和 I/O 系统包含哪些部分。}


\concept{设备管理四个功能}
\begin{points}
  \item 设备分配：根据用户 I/O 请求分配设备、控制器和通道。
  \item 设备处理：启动设备，并在 I/O 完成时处理中断。
  \item 缓冲管理：缓和 CPU 与 I/O 速度不匹配，管理缓冲区分配与释放。
  \item 设备独立性：用户编程使用逻辑设备名，不依赖具体物理设备。
\end{points}
\plain{四个功能最好能一一对上：分配管谁用设备，处理管怎么启动和完成，缓冲管速度差，独立性管逻辑名和物理名解耦。}
\exam{常考设备管理四项基本功能的对应含义。}


\concept{按使用特性分类}
\begin{points}
  \item 存储设备：保存信息，如磁盘、磁带、SSD。
  \item 传输设备：传输数据，如网卡、调制解调器。
  \item 人机交互设备：输入或输出用户信息，如键盘、显示器、打印机。
\end{points}
\plain{按使用特性分类看的是“设备拿来干什么”，不是看它能不能共享或按什么粒度传数据。}
\exam{常考存储设备、传输设备、人机交互设备的代表例子。}


\concept{按共享属性分类}
\begin{points}
  \item 独占设备：一段时间内只允许一个进程使用，如传统打印机。
  \item 共享设备：允许多个进程交替或并发使用，如磁盘。
  \item 虚拟设备：通过虚拟技术把独占设备改造成多个逻辑设备，如 SPOOLing 打印机。
\end{points}
\plain{这组分类要抓住“同一时段能否给多人用”：独占设备一次只给一个，虚拟设备则是把独占设备在逻辑上改成可共享。}
\exam{常考独占、共享、虚拟设备的例子，尤其打印机和磁盘。}


\concept{按信息交换单位分类}
\begin{points}
  \item 块设备以固定大小数据块为单位传输，如磁盘、磁带。
  \item 字符设备以字符流为单位传输，如键盘、鼠标、串口、打印机。
  \item 选择题中，磁盘通常是块设备；键盘、打印机、显示器通常是字符设备。
\end{points}
\plain{块设备和字符设备最容易从“能否随机按块定位”和“传输单位是否固定”来区分。}
\exam{常考磁盘、键盘、打印机分别属于哪类设备。}


\section{I/O 硬件}
\concept{设备控制器与寄存器}
\begin{points}
  \item 设备控制器是 CPU 与外设之间的电子部件，负责控制设备操作。
  \item 控制器通常包含数据寄存器、状态寄存器、控制寄存器。
  \item CPU 通过读写控制器寄存器与设备交互。
  \item 设备完成操作或发生错误时，通过中断通知 CPU。
\end{points}
\plain{设备控制器不是外设本体，而是 CPU 与外设之间的控制桥梁，OS 真正常接触的是这些寄存器。}
\exam{常考控制寄存器、状态寄存器、数据寄存器各自的作用。}


\concept{端口、总线与控制器接口}
\begin{points}
  \item 端口 port 是设备与计算机连接和通信的接口点。
  \item 总线 bus 是多设备共享的一组通信线路及其协议。
  \item 设备控制器内部通常既有与 CPU 的接口，也有与具体设备的接口，以及负责译码和控制的 I/O 逻辑。
\end{points}
\plain{可以把控制器理解成“翻译官 + 调度员”：一边听 CPU 的，一边管具体设备怎么动。}
\exam{常考控制器由哪几部分组成，或端口/总线/控制器分别是什么。}


\concept{I/O 地址空间}
\begin{points}
  \item 端口映射 I/O：I/O 设备有独立地址空间，CPU 使用专门 I/O 指令访问端口。
  \item 内存映射 I/O：设备寄存器映射到内存地址空间，CPU 用普通访存指令访问。
  \item 二者都需要保护，普通用户程序不能随意访问设备寄存器。
\end{points}
\plain{I/O 地址空间这题只看一点：设备寄存器到底是单独编址，还是映射进主存地址空间。}
\exam{常考端口映射 I/O 和内存映射 I/O 的区别。}


\concept{独立编址与内存映像编址比较}
\begin{points}
  \item 独立编址把 I/O 端口地址空间与内存地址空间分开，通常使用专用 I/O 指令。
  \item 内存映像编址把设备寄存器映射到主存地址空间，使用普通访存指令即可访问。
  \item 前者硬件隔离清晰；后者编程更统一，很多体系结构广泛采用。
\end{points}
\plain{一个是“设备地址单开一本账”，一个是“设备也算在内存地址里”。}
\exam{常考两种编址方式的区别，以及哪一种需要专用 IN/OUT 一类指令。}


\concept{通道}
\begin{points}
  \item 通道是一种专门负责 I/O 控制的处理部件，比普通设备控制器更强。
  \item 通道能执行通道程序，组织多个 I/O 操作，减轻 CPU 负担。
  \item 处理机和通道可以并行工作，因此通道提高了大型系统 I/O 并行能力。
\end{points}
\plain{通道比一般控制器更像一个小型 I/O 处理机，能自己按通道程序组织一串 I/O 操作。}
\exam{常考通道和 DMA 的区别，尤其“谁能执行更复杂的控制逻辑”。}


\section{I/O 控制方式}
\concept{程序直接控制方式}
\begin{points}
  \item CPU 反复查询设备状态寄存器，直到设备就绪，再进行数据传送。
  \item 优点：实现简单。
  \item 缺点：CPU 忙等，浪费处理机时间。
  \item 又称轮询方式，适合简单或低速设备。
\end{points}
\plain{程序直接控制就是 CPU 亲自盯着设备状态反复问，简单但最浪费 CPU。}
\exam{常考轮询方式为什么效率低，以及它和中断方式的根本差别。}


\concept{中断驱动方式}
\begin{points}
  \item CPU 启动 I/O 后转去执行其他任务，设备完成后发中断通知 CPU。
  \item 优点：减少 CPU 忙等，提高并发性。
  \item 缺点：每次传送仍需 CPU 参与，频繁中断会增加开销。
\end{points}
\plain{这里要抓住 I/O 主线：CPU 发出 I/O 请求后不用原地等，设备完成时再用中断把 CPU 叫回来收尾。}
\exam{常考中断处理步骤、中断与异常/系统调用区别、为什么中断能提高 CPU 与 I/O 并行度。}


\concept{DMA 方式}
\begin{points}
  \item DMA 允许设备控制器在主存和设备之间直接传送一批数据，无需 CPU 逐字节搬运。
  \item CPU 只负责初始化 DMA 控制器，设置内存地址、传输长度、方向等。
  \item 传输完成后 DMA 控制器发中断通知 CPU。
  \item 优点：大幅减少 CPU 参与，适合磁盘、网卡等高速块设备。
\end{points}
\plain{DMA 的关键不是“没有中断”，而是“搬数据这件事不用 CPU 一字节一字节亲手做”，CPU 只在首尾参与。}
\exam{常考 DMA 和中断方式的差别，尤其“中断发生在每次传送后还是一批传送结束后”。}


\concept{通道控制方式}
\begin{points}
  \item CPU 把一组 I/O 操作写成通道程序交给通道执行。
  \item 通道独立控制设备完成复杂 I/O 序列，完成后中断 CPU。
  \item 与 DMA 相比，通道更像专用 I/O 处理机，控制能力更强。
\end{points}
\plain{通道比 DMA 更进一步：DMA 主要负责一批数据搬运，通道还能按通道程序组织多步 I/O 控制逻辑。}
\exam{比较题常问 DMA 和通道谁更像专用处理机、谁能执行更复杂的 I/O 控制序列。}


\concept{轮询、中断、DMA、通道的核心区别}
\begin{points}
  \item 轮询方式中，CPU 主动反复查状态，忙等最明显。
  \item 中断方式中，每次数据传送仍通常需要 CPU 参与，但不用持续忙等。
  \item DMA 方式中，CPU 只做初始化，一批数据由 DMA 控制器直接在设备和内存间搬运。
  \item 通道方式比 DMA 更进一步，能执行通道程序，控制一组 I/O 操作序列。
\end{points}
\plain{四种方式的主线就是：CPU 参与程度越来越低，I/O 并行能力越来越强。}
\exam{常考比较题：谁在每个字节/每批数据完成后中断 CPU，谁负责真正搬数据。}


\section{I/O 软件原理}
\concept{I/O 软件的设计目标和原则}
\begin{points}
  \item I/O 软件的总体设计目标主要是\key{高效率}和\key{通用性}。
  \item 设计时要重点考虑设备无关性，即上层接口尽量不依赖具体硬件细节。
  \item 还要考虑出错处理，因为 I/O 故障既可能来自设备本身，也可能来自更高层数据结构或访问环境。
  \item 同步传输、异步传输，以及独占设备和共享设备的不同管理方式，也都属于 I/O 软件设计时必须统一处理的问题。
\end{points}
\plain{这块是整章 I/O 软件的总纲：目标不是单纯“能驱动设备”，而是既要高效，又要让不同设备在统一框架下工作。}
\exam{简答题若问 I/O 软件设计原则，标准答法通常先写“高效率、通用性”，再补“设备无关、出错处理、同步/异步、独占/共享设备管理”。}

\concept{I/O 软件层次}
\begin{points}
  \item 用户层 I/O 软件：库函数、假脱机程序等，为用户提供高级接口。
  \item 设备独立性软件：提供统一逻辑设备接口、命名、保护、缓冲、错误处理。
  \item 设备驱动程序：针对具体设备控制器发命令、处理中断。
  \item 中断处理程序：响应设备中断，完成状态检查、唤醒等待进程等。
\end{points}
\plain{I/O 软件分层的主线是“越往上越面向用户和统一接口，越往下越贴近具体硬件和寄存器”。}
\exam{常考四层软件各自负责什么，尤其设备无关软件和驱动程序的边界。}


\concept{I/O 中断处理程序}
\begin{points}
  \item I/O 中断处理程序在设备完成操作、发生错误或出现其他异步事件时被触发。
  \item 主要工作包括检查中断原因、读取状态、确认本次 I/O 是否正常结束，并唤醒等待的驱动程序或进程。
  \item 若发现异常，还要把错误信息上报给更高层软件，交由驱动或设备无关软件进一步处理。
  \item 它位于 I/O 软件的最底层之一，直接承接硬件中断信号。
\end{points}
\plain{驱动程序负责“发命令并组织 I/O”，中断处理程序负责“设备回话后先接住信号并做第一轮收尾”。}
\exam{PPT 单独列了它的功能，答层次题时不要只写“处理中断”，最好写出正常结束、异常事件、唤醒等待者这几类职责。}


\concept{设备驱动程序的职责}
\begin{points}
  \item 接收来自上层的抽象 I/O 请求。
  \item 把逻辑请求转换为设备控制器能理解的具体命令。
  \item 负责设备名到端口地址、逻辑记录到物理记录、逻辑操作到物理操作的转换。
  \item 启动设备，必要时把请求挂入等待队列，并响应控制器发来的中断。
\end{points}
\plain{驱动程序就是“上层说人话、下层说设备话”，它负责把两边翻译通。}
\exam{常考驱动程序和设备无关软件的分工，尤其是谁做地址换算、谁做通用缓冲和保护。}


\concept{设备无关软件的功能细分}
\begin{points}
  \item 设备命名：把逻辑设备名映射到实际设备和驱动程序。
  \item 设备保护：防止未授权进程直接使用设备。
  \item 提供统一逻辑块：屏蔽物理扇区大小差异。
  \item 缓冲管理、存储设备空间管理、独占设备分配与释放、通用错误处理。
\end{points}
\plain{设备无关软件做的是“所有设备都差不多要做的通用活”，不是跟某块网卡或某个磁盘死绑定。}
\exam{常考哪些功能属于设备无关软件，哪些必须由驱动程序完成。}


\concept{阻塞 I/O、非阻塞 I/O 与异步 I/O}
\begin{points}
  \item 阻塞 I/O：调用后进程挂起，直到 I/O 完成。
  \item 非阻塞 I/O：调用立即返回，可能只返回部分数据或暂时返回不能读写。
  \item 异步 I/O：I/O 在后台继续执行，完成后由 I/O 子系统再通知进程。
\end{points}
\plain{阻塞是“你等着”；非阻塞是“你先回去看看别的事”；异步是“事我替你办，办完再通知你”。}
\exam{常考非阻塞和异步的区别，这两个词不能混。}


\concept{设备独立性}
\begin{points}
  \item 设备独立性又称设备无关性，指用户程序不依赖具体物理设备。
  \item 第一层含义：独立于同类设备的具体设备号，例如打印到逻辑打印机而非某台具体打印机。
  \item 第二层含义：在一定抽象下独立于设备类型，例如统一用文件接口访问普通文件和部分设备文件。
  \item 好处：提高可移植性、灵活性和设备分配效率。
\end{points}
\plain{设备独立性的本质是“用户只说要什么设备能力，不必绑定某台具体硬件”，具体映射交给系统做。}
\exam{常考两层含义：独立于同类设备号、以及在一定抽象下独立于设备类型。}


\concept{逻辑设备表 LUT}
\begin{points}
  \item 逻辑设备表用于把逻辑设备名映射为物理设备名和驱动程序入口地址。
  \item 用户程序用逻辑设备名提出请求，系统运行时再完成映射。
  \item 这使 I/O 重定向和设备替换更容易实现。
\end{points}
\plain{LUT 就像“通讯录”：你记逻辑名字，系统去查真正该连哪台设备。}
\exam{常考设备独立性如何实现，标准回答里通常要提逻辑设备表。}


\section{缓冲技术管理}
\concept{为什么需要缓冲}
\begin{points}
  \item CPU 和 I/O 设备速度差异巨大，缓冲可缓和二者速度不匹配。
  \item 缓冲可减少中断次数和设备启动次数，提高并行程度。
  \item 缓冲还能适配不同数据粒度，例如用户按字节读写，磁盘按块传输。
\end{points}
\plain{引入缓冲不只是为了“快慢配速”，还为了把传输粒度、调用语义和并行重叠这几件事都协调起来。}
\exam{老师 PPT 把原因分成三类：速度不匹配、逻辑/物理记录大小不一致、复制语义。}


\concept{引入缓冲的另外两个原因}
\begin{points}
  \item 协调逻辑记录大小与物理记录大小不一致。
  \item 支持复制语义：系统调用返回后，即使用户缓冲区内容改变，也不应影响本次 I/O 应写出的数据。
\end{points}
\plain{缓冲不只是为了解决快慢差，还为了“先把数据稳稳拷住”，避免用户随后改内容把结果搞乱。}
\exam{常考“为什么 write 返回后还能保证写出的是调用当时的数据”，答案通常要提内核缓冲区。}


\concept{单缓冲}
\begin{points}
  \item 系统设置一个缓冲区，设备先把数据送入缓冲区，再由进程取走。
  \item 相比无缓冲，可让设备传输和进程处理部分重叠。
  \item 但一个缓冲区仍可能使设备或进程等待。
\end{points}
\plain{单缓冲下虽然能有部分重叠，但设备输入、系统传送、进程处理仍会围绕同一个缓冲区相互制约。}
\exam{常考单缓冲的时间分析，通常要判断三段时间哪些能并行、哪些必须串行。}


\concept{双缓冲}
\begin{points}
  \item 设置两个缓冲区，一个用于设备填充，另一个供进程处理。
  \item 可以更好地实现输入/输出与处理并行。
  \item 常用于数据流连续、生产和消费速度接近的场景。
\end{points}
\plain{双缓冲的核心直觉就是“一个在装，一个在用”，这样设备和进程更不容易互相等。}
\exam{常考它为什么比单缓冲并行性更强，以及大量数据处理时一块数据的时间怎么分析。}


\concept{循环缓冲}
\begin{points}
  \item 当双缓冲仍不足以平衡 I/O 与处理速度差异时，可继续增加缓冲区数量，构成首尾相连的环形缓冲队列。
  \item 循环缓冲特别适合连续数据流场景，使设备输入、系统传送和进程处理能在更多缓冲区之间流水重叠。
  \item 它更像为一类合作进程定制的连续传输通道。
\end{points}
\plain{循环缓冲的核心不是“缓冲区多”，而是把多个缓冲区按环组织起来，让数据流能一块接一块不断往前推。}
\exam{题目若问什么时候要从双缓冲进一步升级到循环缓冲，关键答“处理大量连续数据、速度差仍较大时”。}

\concept{缓冲池}
\begin{points}
  \item 缓冲池由系统统一管理一组可复用缓冲区，按需分配、回收并供多个进程共享。
  \item 与更偏专用传输通道的循环缓冲不同，缓冲池更适合多设备、多进程同时竞争缓冲资源的场景。
  \item 它既可用于输入，也可用于输出，目标是提高缓冲区整体利用率。
\end{points}
\plain{缓冲池更像操作系统统一维护的一批“公共周转箱”，谁需要谁来借，用完再还。}
\exam{选择题常比较循环缓冲和缓冲池：前者偏连续流专用，后者偏系统范围共享。}


\concept{缓冲区高速缓存}
\begin{points}
  \item 缓冲区高速缓存通常指把常用磁盘块副本保留在内存中，以减少再次访问慢速外设的次数。
  \item 它更接近 cache 的思路，主要利用局部性，而不是单纯协调传输速度。
  \item 对频繁访问的目录块、文件数据块、索引块等，命中后可直接从内存返回，提高 I/O 效率。
  \item 它与单缓冲、双缓冲、循环缓冲并不处于同一比较维度：前者偏“缓存热点块”，后者偏“协调传输过程”。
\end{points}
\plain{这部分容易和普通缓冲混淆，最稳的区分方式是看目标：缓冲区高速缓存是为了“少访盘”，普通缓冲更多是为了“顺利搬数据”。}
\exam{若题目专门问“缓冲”和“缓冲区高速缓存”关系，作答时要写出一个偏传输协调、一个偏局部性复用。}


\concept{缓冲与缓存区别}
\begin{points}
  \item 缓冲 buffer 主要解决速度不匹配和数据传输粒度不一致。
  \item 缓存 cache 主要利用局部性保存数据副本，减少慢速设备访问。
  \item 例子：磁盘块缓存是 cache；键盘输入缓冲区更偏 buffer。
\end{points}
\plain{buffer 更偏“传输时先放一放、调一调节奏”，cache 更偏“常用数据留下副本，下次少访问慢设备”。}
\exam{常考 buffer 和 cache 的目标区别，而不是只让你翻译两个英文单词。}


\concept{缓冲题的时间判断}
\begin{points}
  \item 若从设备输入一块数据到缓冲区时间为 $T$，从缓冲区送到用户区时间为 $M$，用户处理时间为 $C$。
  \item 处理大量数据时，单缓冲下每块数据的处理时间常按 $\max(T,C)+M$ 理解。
  \item 双缓冲下并行性更强，常按 $\max(T,M+C)$ 或结合教材具体流程判断。
  \item 做题时先画出“设备输入、系统传送、进程计算”三段能否重叠，不要直接背公式。
\end{points}
\plain{缓冲计算题本质是在问：这几段时间哪些能重叠，哪些必须串着来。}
\exam{常考单缓冲/双缓冲时间公式或让你自己推，先分清 T、M、C 的含义。}


\section{设备分配与 SPOOLing}
\concept{设备分配中的数据结构}
\begin{points}
  \item 设备控制表 DCT：描述一台设备的状态、类型、队列等。
  \item 控制器控制表 COCT：描述设备控制器状态。
  \item 通道控制表 CHCT：描述通道状态。
  \item 系统设备表 SDT：记录系统中所有设备及其控制表入口。
\end{points}
\plain{这些表的共同作用是把“设备、控制器、通道、等待队列”的状态逐层组织起来，便于系统按层分配资源。}
\exam{常考 DCT、COCT、CHCT、SDT 各自记录什么，以及它们之间的指针连接关系。}


\concept{设备分配过程}
\begin{points}
  \item 根据用户请求的逻辑设备名查找系统设备表。
  \item 分配设备、控制器和通道，若某一级不可用则进程等待。
  \item 分配成功后启动 I/O 操作，I/O 完成后释放相关资源并唤醒等待进程。
\end{points}
\plain{写过程题时按“查逻辑名、分配设备、再分配控制器和通道、启动 I/O、完成后释放”这个顺序展开。}
\exam{常考设备分配过程的步骤题，或让你说明某一级资源不足时进程为什么要等待。}


\concept{设备分配策略}
\begin{points}
  \item 设备分配要考虑设备使用性质、分配算法、安全性和设备独立性。
  \item 常见分配算法有先来先服务和优先级高者优先。
  \item 独占设备、共享设备、虚拟设备的分配方式并不相同。
\end{points}
\plain{设备分配不是“有空就给”，还要看设备类型、请求顺序和是否会带来死锁。}
\exam{常考设备分配时应考虑哪些因素，或区分独享分配、共享分配、虚拟分配。}

\concept{设备分配算法}
\begin{points}
  \item \term{先来先服务}：按进程提出设备请求的先后次序排队，设备优先分配给请求队列队首进程。
  \item \term{优先级高者优先}：先按进程优先级安排设备请求；若优先级相同，再按先来先服务排队。
  \item 这两类算法解决的都是“多个进程竞争同一设备时，系统先满足谁”的问题。
  \item 前者实现简单且相对公平；后者更利于保障关键任务，但可能延长低优先级进程等待时间。
\end{points}
\plain{这块不要和磁盘调度算法混淆：设备分配算法管的是“谁先拿到设备”，磁盘调度算法管的是“磁头接下来怎么走”。}
\exam{选择题常考设备分配算法有哪些，或问“同优先级请求接下来按什么规则排队”。}


\concept{独享分配、共享分配与虚拟分配}
\begin{points}
  \item \term{独享分配}：设备一旦分给某进程，就一直由该进程独占，直到完成或释放。
  \item \term{共享分配}：多个进程可共同使用同一设备，但系统必须安排合理的访问次序。
  \item \term{虚拟分配}：先给进程分配共享存储空间，再由系统在适当时机把这部分数据与真实独占设备交换。
  \item SPOOLing 就是把独占设备改造成逻辑共享设备的典型虚拟分配方法。
\end{points}
\plain{三者的根本区别在于：进程拿到的是“整台真实设备”、还是“被调度地共用设备”、还是“先拿一块共享存储区由系统代办 I/O”。}
\exam{PPT 把这三类分配方式分开讲了，常考打印机为什么更适合虚拟分配而不是直接共享分配。}


\concept{静态分配与动态分配}
\begin{points}
  \item 静态分配在作业开始前一次分配所需全部设备、控制器和通道，运行期间一直占有。
  \item 优点是不会发生设备分配死锁；缺点是利用率低。
  \item 动态分配在进程运行过程中按需分配、用完即释放。
  \item 优点是利用率高，但策略不当可能引发死锁。
\end{points}
\plain{静态分配像“开工前一次领齐工具”，动态分配像“用到什么再去借什么”。}
\exam{常考为什么静态分配安全但效率低，动态分配高效但有死锁风险。}


\concept{安全分配与不安全分配}
\begin{points}
  \item 安全分配下，进程发出 I/O 请求后进入阻塞，直到本次 I/O 完成并释放相关资源。
  \item 不安全分配下，进程发出一个 I/O 请求后还能继续执行，并继续发出后续 I/O 请求。
  \item 前者更安全但推进慢，后者并发性更强但可能造成死锁。
\end{points}
\plain{安全分配像“这件事办完再办下一件”，不安全分配像“先占着这个再去排下一个队”。}
\exam{常考设备分配中的安全性，和死锁章节形成呼应。}


\concept{单通路与多通路分配}
\begin{points}
  \item 单通路系统中，进程常直接按物理设备名申请设备，灵活性较差。
  \item 多通路系统中，用户更适合按逻辑设备名申请，系统再在多台同类设备、控制器、通道之间选择可用路径。
  \item 多通路结构提高了灵活性、可靠性和资源利用率。
\end{points}
\plain{单通路像只能指定某一台机器，多通路则像系统帮你自动找同类里空闲的那一台。}
\exam{常考为什么多通路系统更适合和设备独立性结合。}


\concept{单通路设备分配步骤}
\begin{points}
  \item 单通路分配通常假定进程按物理设备名提出请求，系统按“设备 -> 控制器 -> 通道”顺序逐级分配。
  \item 先由系统设备表 SDT 找到对应设备控制表 DCT，检查设备是否忙，必要时进入设备等待队列。
  \item 再由 DCT 找到控制器控制表 COCT，检查控制器状态；若忙则进入控制器等待队列。
  \item 最后由 COCT 找到通道控制表 CHCT，检查通道状态；若可用则完成分配并启动 I/O。
  \item 任一级资源忙碌，都可能导致进程在该级等待，而不是直接开始 I/O。
\end{points}
\plain{单通路设备分配题最稳的写法就是顺着 `SDT -> DCT -> COCT -> CHCT` 往下找，逐级判断谁忙谁闲。}
\exam{过程题常考每一级查什么状态、忙时进哪个等待队列，以及为什么必须先分到设备再分控制器和通道。}


\concept{SPOOLing 系统}
\begin{points}
  \item SPOOLing 是外部设备联机并行操作技术，常用于把独占设备改造成共享设备。
  \item 基本思想：把慢速独占设备的输入/输出先放到磁盘上的输入井/输出井中，由后台进程统一和实际设备交互。
  \item 典型做法可概括为：\key{预输入}、\key{缓输出}和\key{井管理程序}协同工作。
  \item 系统通常还设置内存中的输入缓冲区和输出缓冲区，用于暂存设备与井之间搬运的数据。
  \item 打印机例子：用户进程把打印数据写入磁盘输出井后即可继续执行，打印进程再按队列送往物理打印机。
  \item SPOOLing 组成：输入井、输出井、输入缓冲区、输出缓冲区、输入进程、输出进程、井管理程序。
  \item 直觉：用户不是直接占用打印机，而是把任务交到“打印队列”；物理打印机仍独占，但逻辑上多个用户可同时提交打印任务。
\end{points}
\plain{SPOOLing 的要害是用磁盘井加后台进程做中转，所以多个用户提交任务时，不必直接排队占物理打印机。}
\exam{常考 SPOOLing 的组成和工作流程，并说明它为什么能把独占设备改造成逻辑共享设备。}
